% =========================================================
% Trigonometric Functions — Notes
% Chapter: elementary-functions
% Section: trigonometric
% Volume III · Analysis
% =========================================================

\section{Trigonometric Functions}
\label{sec:trigonometric}

% ---------------------------------------------------------
% Toolkit
% ---------------------------------------------------------

\begin{tcolorbox}[title={Toolkit --- Trigonometric Functions}]
\begin{itemize}
  \item $\sin x$ and $\cos x$ are defined by absolutely convergent
    power series; properties are derived analytically, not from
    geometry.
  \item Pythagorean identity: $\sin^2 x + \cos^2 x = 1$. Proof:
    the derivative of the left side is $0$ and its value at $0$
    is $1$.
  \item Addition formulas follow from $e^{i(x+y)} = e^{ix}e^{iy}$.
  \item Fundamental trigonometric limit: $\lim_{x \to 0}\sin x/x = 1$.
  \item Second fundamental limit: $\lim_{x\to 0}(1-\cos x)/x^2 = 1/2$.
  \item Inverse functions $\arcsin$, $\arccos$, $\arctan$ are defined
    by restricting $\sin$, $\cos$, $\tan$ to monotone subdomains.
\end{itemize}
\end{tcolorbox}

% ---------------------------------------------------------
% Definition — Sine and Cosine
% ---------------------------------------------------------

\begin{tcolorbox}[title={Definition --- Sine and Cosine}]
\begin{definition}[Sine and Cosine via Power Series]
\label{def:sine-cosine}
The \textbf{sine} and \textbf{cosine} functions are defined for all
$x \in \mathbb{R}$ by the absolutely convergent power series
\[
  \sin x \;:=\;
  \sum_{n=0}^{\infty} \frac{(-1)^n x^{2n+1}}{(2n+1)!}
  \;=\; x - \frac{x^3}{3!} + \frac{x^5}{5!} - \cdots
\]
\[
  \cos x \;:=\;
  \sum_{n=0}^{\infty} \frac{(-1)^n x^{2n}}{(2n)!}
  \;=\; 1 - \frac{x^2}{2!} + \frac{x^4}{4!} - \cdots
\]
\end{definition}
\end{tcolorbox}

\begin{remark*}[Standard quantified statement]
\[
  \forall x \in \mathbb{R} :\;
  \sin x = \sum_{n=0}^\infty \frac{(-1)^n x^{2n+1}}{(2n+1)!},
  \quad
  \cos x = \sum_{n=0}^\infty \frac{(-1)^n x^{2n}}{(2n)!},
\]
both absolutely convergent by the ratio test.
\end{remark*}

\begin{remark*}[Definition predicate reading]
\[
  \operatorname{Sine}(x,y)
  \;\colonequiv\;
  y = \sum_{n=0}^\infty \frac{(-1)^n x^{2n+1}}{(2n+1)!}.
  \qquad
  \operatorname{Cosine}(x,y)
  \;\colonequiv\;
  y = \sum_{n=0}^\infty \frac{(-1)^n x^{2n}}{(2n)!}.
\]
\end{remark*}

\begin{remark*}[Key properties]
\begin{enumerate}[label=(\roman*)]
  \item $\sin 0 = 0$;\; $\cos 0 = 1$.
  \item $\sin(-x) = -\sin x$ (odd);\; $\cos(-x) = \cos x$ (even).
  \item $(\sin x)' = \cos x$;\; $(\cos x)' = -\sin x$.
  \item $|\sin x| \leq 1$;\; $|\cos x| \leq 1$ for all $x$.
  \item $\sin^2 x + \cos^2 x = 1$ (Pythagorean identity).
  \item Periodic: $\sin(x+2\pi) = \sin x$;\; $\cos(x+2\pi) = \cos x$.
\end{enumerate}
\end{remark*}

\begin{remark*}[Relationship to $e^{ix}$]
Euler's formula: $e^{ix} = \cos x + i\sin x$, where $e^{ix}$ is
defined by the complex exponential series. This gives $\cos x =
\operatorname{Re}(e^{ix})$ and $\sin x = \operatorname{Im}(e^{ix})$.
Euler's identity $e^{i\pi} + 1 = 0$ follows immediately.
\end{remark*}

% ---------------------------------------------------------
% Proposition — Pythagorean Identity
% ---------------------------------------------------------

\begin{proposition}[Pythagorean Identity]
\label{prop:pythagorean-identity}
For all $x \in \mathbb{R}$:
\[
  \sin^2 x + \cos^2 x = 1.
\]

\smallskip\noindent
\hyperref[prf:pythagorean-identity]{\textit{Go to proof.}}
\end{proposition}

\begin{remark*}[Standard quantified statement]
\[
  \forall x \in \mathbb{R} :\; \sin^2 x + \cos^2 x = 1.
\]
\end{remark*}

\begin{remark*}[Negated quantified statement]
\[
  \neg\,(\forall x:\sin^2 x+\cos^2 x=1)
  \;\iff\;
  \exists x\in\mathbb{R}:\sin^2 x+\cos^2 x\neq 1.
\]
The negation has no witness --- the identity holds universally.
\end{remark*}

\begin{remark*}[Definition predicate reading]
\[
  \operatorname{PythagoreanIdentity}(x)
  \;\colonequiv\;
  \operatorname{Sine}(x,s) \;\wedge\; \operatorname{Cosine}(x,c)
  \;\wedge\; s^2 + c^2 = 1.
\]
\end{remark*}

\begin{remark*}[Proof strategy]
Define $f(x) := \sin^2 x + \cos^2 x$. Then $f'(x) = 2\sin x \cos x
+ 2\cos x(-\sin x) = 0$, so $f$ is constant. Since $f(0) = 0 + 1 =
1$, the identity follows.
\end{remark*}

% ---------------------------------------------------------
% Proposition — Addition Formulas
% ---------------------------------------------------------

\begin{proposition}[Addition Formulas]
\label{prop:trig-addition-formulas}
For all $x, y \in \mathbb{R}$:
\[
  \sin(x+y) = \sin x \cos y + \cos x \sin y,
  \qquad
  \cos(x+y) = \cos x \cos y - \sin x \sin y.
\]

\smallskip\noindent
\hyperref[prf:trig-addition-formulas]{\textit{Go to proof.}}
\end{proposition}

\begin{remark*}[Standard quantified statement]
\[
  \forall x,y\in\mathbb{R}:\;
  \sin(x+y) = \sin x\cos y + \cos x\sin y
  \;\wedge\;
  \cos(x+y) = \cos x\cos y - \sin x\sin y.
\]
\end{remark*}

\begin{remark*}[Negated quantified statement]
\[
  \neg\,(\text{addition formula for }\sin)
  \;\iff\;
  \exists x,y\in\mathbb{R}:\;
  \sin(x+y) \neq \sin x\cos y + \cos x\sin y.
\]
The negation has no witness.
\end{remark*}

\begin{remark*}[Definition predicate reading]
\[
  \operatorname{TrigAddition}(x,y)
  \;\colonequiv\;
  \sin(x+y) = \sin x\cos y + \cos x\sin y
  \;\wedge\;
  \cos(x+y) = \cos x\cos y - \sin x\sin y.
\]
\end{remark*}

\begin{remark*}[Proof strategy]
From the Cauchy product: $e^{i(x+y)} = e^{ix}e^{iy}$. Taking real
and imaginary parts yields both addition formulas simultaneously.
\end{remark*}

\begin{remark*}[Derived identities]
\begin{align*}
  \sin(2x) &= 2\sin x \cos x \\
  \cos(2x) &= \cos^2 x - \sin^2 x = 1 - 2\sin^2 x = 2\cos^2 x - 1 \\
  \sin^2 x &= \tfrac{1}{2}(1 - \cos 2x) \\
  \cos^2 x &= \tfrac{1}{2}(1 + \cos 2x)
\end{align*}
\end{remark*}

% ---------------------------------------------------------
% Definition — Tangent and other trig functions
% ---------------------------------------------------------

\begin{definition}[Tangent, Cotangent, Secant, Cosecant]
\label{def:other-trig}
\[
  \tan x := \frac{\sin x}{\cos x},\quad
  \cot x := \frac{\cos x}{\sin x},\quad
  \sec x := \frac{1}{\cos x},\quad
  \csc x := \frac{1}{\sin x}.
\]
$\tan x$ and $\sec x$ are defined where $\cos x \neq 0$,
i.e., $x \neq \pi/2 + k\pi$ for $k \in \mathbb{Z}$.
\end{definition}

\begin{remark*}[Standard quantified statement]
\[
  \forall x\in\mathbb{R},\;\cos x\neq 0:\;
  \tan x = \frac{\sin x}{\cos x}.
\]
\end{remark*}

\begin{remark*}[Definition predicate reading]
\[
  \operatorname{Tangent}(x,y)
  \;\colonequiv\;
  \cos x \neq 0 \;\wedge\; y = \frac{\sin x}{\cos x}.
\]
\end{remark*}

\begin{remark*}[Pythagorean identities for $\tan$]
Dividing $\sin^2 x + \cos^2 x = 1$ by $\cos^2 x$:
$1 + \tan^2 x = \sec^2 x$.
Dividing by $\sin^2 x$: $\cot^2 x + 1 = \csc^2 x$.
\end{remark*}

% ---------------------------------------------------------
% Definition — Inverse Trigonometric Functions
% ---------------------------------------------------------

\begin{definition}[Inverse Trigonometric Functions]
\label{def:inverse-trig}
\begin{itemize}
  \item $\arcsin : [-1,1] \to [-\pi/2, \pi/2]$ is the inverse of
    $\sin\big|_{[-\pi/2,\pi/2]}$.
  \item $\arccos : [-1,1] \to [0,\pi]$ is the inverse of
    $\cos\big|_{[0,\pi]}$.
  \item $\arctan : \mathbb{R} \to (-\pi/2, \pi/2)$ is the inverse
    of $\tan\big|_{(-\pi/2,\pi/2)}$.
\end{itemize}
\end{definition}

\begin{remark*}[Standard quantified statement]
\[
  \forall y\in[-1,1]:\;\sin(\arcsin y)=y.
  \qquad
  \forall x\in[-\pi/2,\pi/2]:\;\arcsin(\sin x)=x.
\]
\end{remark*}

\begin{remark*}[Definition predicate reading]
\[
  \operatorname{Arctan}(y,x)
  \;\colonequiv\;
  x\in(-\pi/2,\pi/2) \;\wedge\; \tan x = y.
\]
\end{remark*}

\begin{remark*}[Negated quantified statement]
\[
  \neg\,\operatorname{Arctan}(y,x)
  \;\iff\;
  x\notin(-\pi/2,\pi/2) \;\vee\; \tan x \neq y.
\]
\end{remark*}

\begin{remark*}[Key limits of $\arctan$]
\[
  \lim_{x \to +\infty} \arctan x = \frac{\pi}{2},
  \qquad
  \lim_{x \to -\infty} \arctan x = -\frac{\pi}{2}.
\]
$\arctan$ is bounded, strictly increasing, and continuous on
$\mathbb{R}$; the Monotone Limit Theorem gives the limits as the
supremum and infimum of its range.
\end{remark*}

% ---------------------------------------------------------
% Theorem — Fundamental Trigonometric Limit
% ---------------------------------------------------------

\begin{theorem}[Fundamental Trigonometric Limit]
\label{thm:fundamental-trig-limit}
\[
  \lim_{x \to 0} \frac{\sin x}{x} = 1.
\]

\smallskip\noindent
\hyperref[prf:fundamental-trig-limit]{\textit{Go to proof.}}
\end{theorem}

\begin{remark*}[Standard quantified statement]
\[
  \forall\varepsilon > 0,\;
  \exists\delta > 0 :\;
  0 < |x| < \delta \Rightarrow
  \left|\frac{\sin x}{x} - 1\right| < \varepsilon.
\]
\end{remark*}

\begin{remark*}[Negated quantified statement]
\[
  \neg\left(\lim_{x\to 0}\frac{\sin x}{x}=1\right)
  \;\iff\;
  \exists\varepsilon_0>0,\;\forall\delta>0,\;\exists x:\;
  0<|x|<\delta \;\wedge\;
  \left|\frac{\sin x}{x}-1\right|\geq\varepsilon_0.
\]
\end{remark*}

\begin{remark*}[Definition predicate reading]
\[
  \operatorname{FunctionLimit}\!\left(\frac{\sin x}{x},\,0,\,1\right)
  \;\colonequiv\;
  \forall\varepsilon>0,\;\exists\delta>0,\;\forall x:\;
  \operatorname{InputTolerance}(x,0,\delta)
  \Rightarrow
  \operatorname{OutputTolerance}\!\left(\frac{\sin x}{x},1,\varepsilon\right).
\]
\end{remark*}

\begin{remark*}[Proof strategy --- series]
From the series: $\sin x / x = 1 - x^2/3! + x^4/5! - \cdots \to 1$
as $x \to 0$ by continuity of power series at their center.
\end{remark*}

\begin{remark*}[Proof strategy --- geometric squeeze]
For $0 < x < \pi/2$, area comparison of two triangles and a circular
sector gives $\cos x \leq \sin x / x \leq 1$. Since $\cos x \to 1$
as $x \to 0^+$, the Squeeze Theorem gives $\sin x / x \to 1$.
Even symmetry extends this to $x \to 0^-$.
\end{remark*}

\begin{remark*}[Significance]
This is the statement $(\sin x)'\big|_{x=0} = 1$, which implies
$(\sin x)' = \cos x$ everywhere. It underlies all trigonometric
differentiation.
\end{remark*}

% ---------------------------------------------------------
% Proposition — Second Fundamental Trigonometric Limit
% ---------------------------------------------------------

\begin{proposition}[Second Fundamental Trigonometric Limit]
\label{prop:limit-one-minus-cos}
\[
  \lim_{x \to 0} \frac{1 - \cos x}{x^2} = \frac{1}{2}.
\]

\smallskip\noindent
\hyperref[prf:limit-one-minus-cos]{\textit{Go to proof.}}
\end{proposition}

\begin{remark*}[Standard quantified statement]
\[
  \forall\varepsilon>0,\;\exists\delta>0:\;
  0<|x|<\delta \Rightarrow
  \left|\frac{1-\cos x}{x^2}-\frac{1}{2}\right|<\varepsilon.
\]
\end{remark*}

\begin{remark*}[Negated quantified statement]
\[
  \neg\left(\lim_{x\to 0}\frac{1-\cos x}{x^2}=\frac{1}{2}\right)
  \;\iff\;
  \exists\varepsilon_0>0,\;\forall\delta>0,\;\exists x:\;
  0<|x|<\delta \;\wedge\;
  \left|\frac{1-\cos x}{x^2}-\frac{1}{2}\right|\geq\varepsilon_0.
\]
\end{remark*}

\begin{remark*}[Definition predicate reading]
\[
  \operatorname{FunctionLimit}\!\left(\frac{1-\cos x}{x^2},\,0,\,\tfrac{1}{2}\right)
  \;\colonequiv\;
  \forall\varepsilon>0,\;\exists\delta>0,\;\forall x:\;
  \operatorname{InputTolerance}(x,0,\delta)
  \Rightarrow
  \operatorname{OutputTolerance}\!\left(\frac{1-\cos x}{x^2},\tfrac{1}{2},\varepsilon\right).
\]
\end{remark*}

\begin{remark*}[Proof strategy]
Use the half-angle identity $1 - \cos x = 2\sin^2(x/2)$:
\[
  \frac{1 - \cos x}{x^2}
  = \frac{2\sin^2(x/2)}{x^2}
  = \frac{1}{2}\left(\frac{\sin(x/2)}{x/2}\right)^2
  \to \frac{1}{2} \cdot 1^2 = \frac{1}{2}.
\]
\end{remark*}
